\chapter{METHODOLOGY}

\pagebreak

\begin{center}
{\LARGE\textbf{METHODOLOGY}}
\end{center}

This chapter describes the approach used to design, implement, and evaluate the TrapPot system \cite{franco2021survey,rodriguezgomez2025systematic,lanz2025optimizing}. The methodology covers the research design, data collection procedures, preprocessing and feature engineering pipeline, Random Forest classifier training, validation plan, and ethical considerations \cite{lee2021classification,hindy2023comprehensive,karim2024efficient}. The chapter explains how TrapPot captures IoT-targeted activity, converts it into useful features, and evaluates classification performance against the project target of at least 99\% accuracy \cite{lee2021classification,karim2024efficient,kalaria2024iotpredictor}.

% ============================================================
\section{Research Design}
\label{sec:research_design}
% ============================================================

The TrapPot study follows a case study research design with experimental elements, supported by a quantitative approach \cite{franco2021survey,rodriguezgomez2025systematic}. This design fits the project because TrapPot is a built system rather than a survey-only study. The research examines how a specific implementation captures attacker behavior, processes logs, and classifies network traffic under controlled conditions \cite{lee2021classification,lanz2025optimizing}. The quantitative approach is used because system logs, feature vectors, and model metrics can be measured and compared statistically \cite{hindy2023comprehensive,mary2024network}.

% ------------------------------------------------------------
\subsection{Research Design Type: Case Study}
\label{subsec:case_study}
% ------------------------------------------------------------

The primary research design is a case study of one concrete deployment: the TrapPot system running in a dedicated and isolated environment \cite{albaseer2024fedpot,commey2025blockchain}. This design allows detailed observation of how medium-interaction honeypot technology and machine learning behave when combined in one practical IoT security workflow \cite{lee2021classification,9359711}.

By focusing on one deployment, the case study makes it possible to observe \cite{rodriguezgomez2025systematic,lanz2025optimizing}:

\begin{itemize}
    \item how attackers interact with emulated Telnet and SSH services \cite{9071014,9558948,famera2025mirai},
    \item how logs move through the ELK-based data pipeline \cite{9001068,log2evt2025}, and
    \item how the Random Forest classifier responds to brute-force, scanning, malware download, command execution, and normal traffic patterns \cite{lee2021classification,karim2024efficient,saied2023comparative}.
\end{itemize}

This depth of analysis is more suitable for a system-development project than a broad theoretical survey because it evaluates the practical integration of honeypots, network monitoring, and AI-based classification.

% ------------------------------------------------------------
\subsection{Research Approach: Quantitative}
\label{subsec:quantitative}
% ------------------------------------------------------------

The research uses quantitative methods because the collected evidence is machine-generated and structured. Cowrie records authentication attempts, commands, sessions, and file activity, while Zeek records ports, protocols, services, packet counts, byte counts, connection duration, and connection state \cite{9071014,log2evt2025}. These records are converted into feature vectors that can be used for statistical analysis and model evaluation \cite{dimauro2021supervised,ahsan2024effects}.

The methodology follows a system-development and experimental-validation pattern \cite{lee2021classification,rodriguezgomez2025systematic}:

\begin{itemize}
	\item Develop and deploy TrapPot using Cowrie, Zeek, the ELK Stack, and a Random Forest classifier in an isolated network \cite{9359711,commey2025blockchain}.
	\item Collect attack traffic data over the specified operational period \cite{lee2021classification,suareztangil2023stories}.
	\item Train and test the Random Forest model using the prepared dataset \cite{saied2023comparative,karim2024efficient}.
	\item Evaluate the model using accuracy, precision, recall, F1-score, and confusion matrix results \cite{confusion2023pmc,mary2024network}.
\end{itemize}

This approach allows the team to measure both the classification performance of the model and the operational behavior of the complete TrapPot pipeline.


% ============================================================
\section{Data Collection}
\label{sec:data_collection}
% ============================================================

TrapPot collects and uses data in two phases. The first phase uses existing research and public datasets to prepare and validate the AI module. The second phase uses logs generated by the deployed honeypot and monitoring system for live classification and operational evaluation \cite{rodriguezgomez2025systematic,lanz2025optimizing}.

% ------------------------------------------------------------
\subsection{Pre-Implementation Data Collection}
\label{subsec:pre_implementation_data}
% ------------------------------------------------------------

Before live deployment, secondary data from public and academic sources was used to design, train, and validate the AI module \cite{hindy2023comprehensive,dekeersmaeker2023survey,alex2023comprehensive}. This makes the model development stage grounded in real IoT security patterns and established research.

\textbf{Sources of Secondary Data:}

\textbf{Academic Publications \& Research Datasets}

\begin{itemize}
    \item Lee et al. (2021) provided labeled attack logs from Cowrie honeypots in IoT smart factory environments, classifying brute-force, scanning, and malware download attacks with 96\% accuracy using Random Forest \cite{lee2021classification}.
    \item Saputro et al. (2021) supplied performance benchmarks and attack traces from medium-interaction IoT honeypots, confirming Cowrie's low resource overhead below 6.3\% CPU \cite{9359711}.
    \item Saied et al. (2023) offered comparative evidence for IoT botnet classification, demonstrating Random Forest's 99.99\% accuracy for multi-class attack detection \cite{saied2023comparative}.
    \item Alajrami \& Abu-Naser (2025) shared an SSH/Telnet honeypot-based intrusion detection approach that supports testing multi-algorithm IDS performance \cite{alajrami2025aipowered}.
\end{itemize}

\textbf{Public IoT Security Datasets}

\begin{itemize}
    \item Mirai botnet traces provide historical attack data from Mirai and variants such as Gafgyt and Mozi, focusing on credential-based attacks against Telnet and SSH \cite{203628,famera2025mirai}.
    \item Cowrie public logs provide examples of SSH and Telnet honeypot sessions that can support feature extraction and model preparation \cite{9071014,9558948}.
    \item IoT-23 and CIC-IoT-2023 provide labeled IoT network traffic datasets that help validate feature relevance and model generalization \cite{iot23dataset,dekeersmaeker2023survey,alex2023comprehensive,christy2024iot}.
\end{itemize}

\textbf{Technical Documentation \& Cybersecurity Reports}

\begin{itemize}
    \item Cowrie documentation provides JSON log structures, command emulation behavior, and session metadata formats \cite{9071014,log2evt2025}.
    \item ELK Stack guidance supports Logstash pipelines and Elasticsearch mappings for parsing and storing honeypot data \cite{9001068}.
    \item Threat intelligence reports and academic surveys provide context on scanning patterns, malware delivery, and IoT attack trends \cite{mahmood2024exploring,islam2023large}.
\end{itemize}

% ------------------------------------------------------------
\subsection{Post-Implementation Data Collection}
\label{subsec:post_implementation_data}
% ------------------------------------------------------------

After deployment in an isolated network environment, TrapPot collects primary operational data in real time \cite{lee2021classification,suareztangil2023stories}. This data is used to:

\begin{itemize}
	\item feed the live classification module,
	\item support future model refinement and validation, and
	\item provide current threat intelligence for analysis and reporting.
\end{itemize}

The implementation and testing details of this phase are presented in Chapter 4, while the observed results are discussed in Chapter 5.

% ============================================================
\section{Data Analysis}
\label{sec:data_analysis}
% ============================================================

% ------------------------------------------------------------
\subsection{Data Acquisition and Preprocessing Pipeline}
\label{subsec:preprocessing_pipeline}
% ------------------------------------------------------------

The data acquisition pipeline relies on Zeek transport-layer outputs. Cowrie acts as the deception layer, while Zeek monitors network connections and generates the tabular stream used by the AI detector \cite{9071014,log2evt2025}. The ELK Stack, consisting of Elasticsearch, Logstash, and Kibana, handles processing, storage, and visualization \cite{9001068}. Logstash processes raw JSON records and extracts fields such as session identifiers, timestamps, source IP addresses, commands, and credential-attempt information \cite{9001068,log2evt2025}.

\textbf{Data Cleaning and Organization:} To maintain data integrity, records are organized by session ID and event timestamp. This helps identify repeated records and preserve the relationship between attacker actions and network events \cite{rodriguezgomez2025systematic}. Missing command or input fields are handled during feature calculation so incomplete sessions do not produce invalid model input \cite{ahsan2024effects,mary2024network}.

\textbf{Ground Truth Labeling Strategy:} As a supervised learning project, accurate labeling is essential \cite{hopf2021supervised,dimauro2021supervised}. A hierarchical strategy assigns one of five target classes to each session or connection record: Brute Force, Scanning, Malware Download, Command Execution, or Normal traffic \cite{lee2021classification,karim2024efficient}.

\begin{itemize}
    \item Malware Download is identified by the presence of data in the URLs field \cite{9071014}.
    \item Command Execution requires evidence of a successful login followed by executed commands \cite{9071014,9558948}.
    \item Brute Force is labeled when there is a high frequency of failed login attempts \cite{203628,famera2025mirai}.
    \item Scanning is categorized by low-interaction sessions that show port or service probing behavior \cite{zhang2023identification,suareztangil2023stories}.
    \item Normal traffic represents baseline operational events, minor polling, or protocol checks with minimal packet counts and predictable transport signatures \cite{lee2021classification}.
\end{itemize}

% ------------------------------------------------------------
\subsection{Feature Engineering and Selection}
\label{subsec:feature_engineering}
% ------------------------------------------------------------

Feature engineering transforms categorical and sequential log data into a compact set of numerical and categorical features that can separate the five classes \cite{dimauro2021supervised,christy2024iot}. These features must be meaningful to the Random Forest model and understandable to the project team, so the prototype avoids complex dimensionality reduction techniques \cite{lee2021classification,dimauro2021supervised}.

The primary features are based on authentication behavior, session activity, command usage, file transfer indicators, and network connection properties \cite{christy2024iot,karim2024efficient}:

\textbf{Authentication Metrics} include the number of unique login attempts per session. A high value indicates automated brute-force activity that tests many credential pairs \cite{203628,famera2025mirai,9558948}.

\textbf{Session Metrics} include session duration in seconds. Short durations can indicate automated scanning, while longer sessions may indicate interactive command execution \cite{lee2021classification,karim2024efficient}.

\textbf{Command Metrics} focus on post-login behavior, especially the ratio of recognized operating system commands to invalid or empty inputs. This provides evidence of command execution after access is obtained \cite{9071014,dimauro2021supervised}.

\textbf{Malware Metrics} use a binary file-transfer flag to indicate whether an upload or download attempt occurred. This feature supports identification of malware download behavior \cite{9071014,9558948}.

\textbf{Network Metrics} are derived from Zeek \texttt{conn.log}, including source and destination ports, duration, byte counts, packet counts, protocol, service, and connection state. These features match the trained detector pipeline and allow live classification from Zeek records.

This feature set supports the project goal of interpretable and efficient classification while remaining compatible with the implemented detector \cite{lee2021classification,saied2023comparative,karim2024efficient}.

% ------------------------------------------------------------
\subsection{Machine Learning Model Training and Optimization}
\label{subsec:model_training}
% ------------------------------------------------------------

The analysis relies on the Random Forest Classifier, a choice supported by literature reporting high accuracy in IoT intrusion detection and low inference latency \cite{lee2021classification,saied2023comparative,karim2024efficient,christy2024iot}. The prepared dataset is partitioned into an 80\% training set and a 20\% hold-out testing set \cite{optimized2024arxiv,karim2024efficient}. Five-fold cross-validation is used during training to evaluate model robustness across the training data \cite{pan2024xgboost,optimized2024arxiv}. To reduce the effect of class imbalance, the model uses class weighting and balancing strategies so that less frequent classes such as Malware Download are not ignored \cite{9213728,9912910,abdulsalam2023multiclass}. Hyperparameters, including the number of decision trees and maximum tree depth, are tuned to support the project's minimum 99\% classification accuracy goal \cite{pan2024xgboost,optimized2024arxiv}.

% ============================================================
\section{Validation Plan}
\label{sec:validation_plan}
% ============================================================

The validation plan measures both the analytical accuracy of the Random Forest classifier and the operational feasibility of TrapPot as a real-time threat intelligence system \cite{confusion2023pmc,mary2024network}. The main success target is classification accuracy above 99\% \cite{lee2021classification,karim2024efficient,kalaria2024iotpredictor}.

% ------------------------------------------------------------
\subsection{Classification Performance Metrics}
\label{subsec:performance_metrics}
% ------------------------------------------------------------

Accuracy is the primary metric because it directly measures the percentage of correct predictions across the five classes \cite{karim2024efficient,mary2024network}. However, security datasets often contain imbalanced classes, so additional metrics are also needed \cite{confusion2023pmc,mary2024network}.

\textbf{Primary Metric:} Accuracy is used to determine whether the model achieves the 99\% classification target across the five attack categories \cite{karim2024efficient,mary2024network}.

\textbf{Secondary Critical Metric:} F1-score is reported alongside accuracy because it balances precision and recall \cite{confusion2023pmc,karim2024efficient}. For TrapPot, a high F1-score helps confirm that less frequent but important classes, such as Malware Download and Command Execution, are not hidden by more frequent brute-force or normal traffic records \cite{lee2021classification,9213728}.

% ------------------------------------------------------------
\subsection{Validation Dataset and Testing Method}
\label{subsec:validation_dataset}
% ------------------------------------------------------------

Final model validation uses the 20\% hold-out dataset that was separated during data preparation \cite{optimized2024arxiv,karim2024efficient}. This data is not used during training, which gives a fairer measure of model generalization \cite{lee2021classification,saied2023comparative,kalaria2024iotpredictor}.

\textbf{Data Collection Duration:} The study design keeps the 30-day operational data collection period for the Cowrie honeypot in an isolated network environment \cite{lee2021classification,suareztangil2023stories}. This duration is intended to provide a stronger sample of attack traffic targeting Telnet and SSH \cite{famera2025mirai,9558948}.

\textbf{Methodology:} The trained Random Forest model is applied to the hold-out set, and the resulting predictions are compared against the ground truth labels established during preprocessing \cite{confusion2023pmc,mary2024network}.

% ------------------------------------------------------------
\subsection{Operational Success Criteria}
\label{subsec:operational_criteria}
% ------------------------------------------------------------

Beyond analytical metrics, TrapPot must demonstrate that it can operate as a real-time monitoring and classification system \cite{rodriguezgomez2025systematic,lanz2025optimizing}. Two system-level criteria are used:

\textbf{Real-Time Inference Latency:} The trained Random Forest model must classify new connection records with low latency so that Kibana can support near-real-time monitoring through the ELK Stack \cite{9001068,karim2024efficient}.

\textbf{Resource Efficiency:} Cowrie, Zeek, Logstash, Elasticsearch, Kibana, and the Random Forest detector must run with acceptable CPU and memory use for a local academic deployment \cite{9359711,alzahrani2025lightweight}. This validates the system's suitability for educational and non-enterprise lab environments \cite{franco2021survey,morozov2024sweet}.

% ------------------------------------------------------------
\subsection{Results Presentation}
\label{subsec:results_presentation}
% ------------------------------------------------------------

The system results are presented through both quantitative tables and visual monitoring outputs \cite{confusion2023pmc,mary2024network}. A confusion matrix is used to show how predicted labels compare with actual labels across the five classes \cite{confusion2023pmc,karim2024efficient,mary2024network}. Kibana dashboards are used to show operational evidence, including attack activity over time, source IP addresses, protocol distribution, commands captured, and AI detector output \cite{9001068}.

% ============================================================
\section{Research Ethics}
\label{sec:research_ethics}
% ============================================================

TrapPot follows a do-no-harm ethical approach. The project observes malicious behavior for research and learning, but it must not endanger university systems, external networks, or identifiable individuals \cite{commey2025blockchain,franco2021survey}. Because the project does not involve human-subject experiments, the main ethical concerns are network containment, responsible log handling, and careful reporting \cite{franco2021survey,alawadi2022_iot_honeypots}.

% ------------------------------------------------------------
\subsection{Network Isolation and Containment}
\label{subsec:network_isolation}
% ------------------------------------------------------------

Network safety is the main ethical and technical requirement \cite{commey2025blockchain,albaseer2024fedpot}. Cowrie is designed to capture attack data, but the deployment must not become a risk to the university network or a launchpad for third-party attacks \cite{commey2025blockchain,fan2019honeydoc}.

\textbf{Strict Segregation:} TrapPot is deployed inside a dedicated and isolated lab network or virtual network segment, separated from production infrastructure \cite{albaseer2024fedpot,fan2019honeydoc}.

\textbf{Preventing Egress:} Firewall and network controls must restrict outgoing traffic from the honeypot environment so that attackers cannot use the simulated system to attack external targets \cite{commey2025blockchain}.

\textbf{Sandbox Security:} Cowrie's medium-interaction design provides a controlled shell environment. Attacker commands are recorded inside the honeypot and do not execute on the host operating system \cite{9071014,9359711}.

% ------------------------------------------------------------
\subsection{Data Handling and Anonymization}
\label{subsec:data_handling}
% ------------------------------------------------------------

Captured logs may include source IP addresses, timestamps, credentials attempted by attackers, and commands entered during sessions \cite{rodriguezgomez2025systematic,lee2021classification}. The project therefore keeps only the data needed for threat intelligence, classification, and academic reporting.

\textbf{Data Minimization:} TrapPot stores the fields required for analysis, such as source IP address, timestamp, credential attempt, command input, and session metadata \cite{rodriguezgomez2025systematic,lee2021classification}.

\textbf{Storage and Access:} Log data stored in ELK and project machines is restricted to the project team and used only for research and evaluation \cite{9001068}.

\textbf{Anonymity in Reporting:} Public reporting uses aggregated statistics and anonymized summaries. Individual raw source IP addresses are not disclosed in the thesis, presentation, or public materials \cite{franco2021survey,rodriguezgomez2025systematic}.

% ------------------------------------------------------------
\subsection{Regulatory Compliance and Log Management}
\label{subsec:regulatory_compliance}
% ------------------------------------------------------------

\textbf{Institutional Approval:} Any 30-day data collection period or external exposure must be approved by the relevant university or lab authority to confirm that the deployment and isolation plan comply with security policies \cite{franco2021survey,alawadi2022_iot_honeypots}.

\textbf{Data Disposal:} After project completion, raw identifiable logs should be purged from ELK and associated storage. Only anonymized and aggregated summaries should be retained for future comparison or academic reference \cite{rodriguezgomez2025systematic,franco2021survey}.
